\section{Introduction}
\label{sec:introduction}

% ROUGH OUTLINE -- one bullet = one point (roughly one paragraph in the final
% text). For co-author feedback; details added during writing.

\begin{itemize}
  \item \textbf{Motivation.} Urban wind fields drive dispersion, pedestrian
        comfort, and wind loading; the dominant uncertainties are the approach
        flow (\inflowangle, \velmag, \shearexp) and the turbulence closure
        (\sgs), not the numerics \cite{britter2003flow,xie2009modelling}.
  \item \textbf{Inverse problem.} Urban sensor data are sparse and noisy;
        ensemble data assimilation (\esmda) can recover forcing and closure
        parameters by fusing observations with a CFD model, treating the solver
        as a derivative-free black box
        \cite{emerick2013esmda,kikumoto2018bayesian,xue2017inverse}.
  \item \textbf{The gap.} Validation is almost always a \twinexp: synthetic
        truth generated by the \emph{same} solver used for assimilation --- the
        \invcrime\ \cite{kaipio2005statistical,wirgin2004inverse}. This sets
        \discrepancy\ to zero by construction and flatters the method; turbulent
        urban \les\ (closure- and geometry-sensitive) is where this bites
        hardest.
  \item \textbf{Contributions.}
        \begin{itemize}
          \item A \crossmodel\ evaluation protocol: truth and assimilation
                solvers differ, spanning pairings of \lbmmodel, \pyudales,
                \pypalm, so \discrepancy\ is present and controllable.
          \item Evidence that \twinexp\ success does not transfer cross-model.
          \item A systematic study of which parameters (forcing vs.\ closure)
                and which algorithmic choices (\Ne, \Na, \Nw, \Cd, localization,
                sensors) govern cross-model success.
          \item Practical guidance for setting up robust urban-flow assimilation.
        \end{itemize}
  \item \textbf{Paper organization.} One sentence per section.
\end{itemize}
